\documentclass[../main]{subfiles}
\begin{document}

\chapter*{Planificación Ciclo Lectivo \ano}
\section*{Introducción}
\begin{itemize}
    
\item ESPACIO CURRICULAR: \materiaLong
\item CURSO:\curso
\item PROFESOR/A:\nombre
\item ORIENTACIÓN EN: Agro-Pecuaria Pecuaria
\item PERFIL DEL EGRESADO:
\item ACUERDOS TRANSVERSALES INSTITUCIONALES (Fluidez lectora y comprensión, Convivencia escolar, etc.).
\item ACUERDOS DEL ÁREA:
    \begin{itemize}
    \item Desarrollar prácticas educativas a través de los proyectos didácticos productivos.
    \item Expresar por medio de las habilidades y capacidades adquiridas los contenidos conceptuales.
    \item Exponer oralmente  utilizando técnicas de comunicación como cierre de cada eje temático.
    \item Proyectar sobre la formación de técnicos agropecuarios con habilidades y capacidades para la inserción y  transformación del sector agropecuario en la zona de influencia y a estudios superiores.
    \end{itemize}
    \item APORTE ESPECÍFICO DEL ESPACIO CURRICULAR AL PERFIL DEL EGRESADO.
    Este espacio se sustenta en los saberes desarrollados en Prácticas de
Producción Regional, Bases Bioagropecuarias y la Producción Vegetal, así como
Ecología Agraria y Economía Social, y por otra parte, se interrelaciona con Riego y
Drenaje, las Producciones Vegetales y Animales y sus respectivos procesos de
procesamiento e industrialización.

La presente propuesta formativa integra teoría y práctica, conforme lo
establece el artículo 35 de la Resolución CFE 229/14. Cada docente deberá garantizar
que al menos un tercio del total de las horas de enseñanza –semanales se dediquen al
desarrollo de prácticas de distinta índole, razón por la cual se recomienda trabajar con
la metodología del aula-taller. 
\end{itemize}
\newpage
\end{document}
